\scanpage{378}
\mspage{490}

\[
(29)\quad
\left\{\begin{aligned}
\underset{\substack{\Vert\,\mathrm{dfn.}\\\underline{u}\rho\underline{u}^{-1}}}{\underline{u}(\rho)}
 &=\underset{\substack{\Vert\\\operatorname{int}(\alpha)\rho}}{\underbrace{\xi\rho}},
 &\underline{u}(\sigma)&=\underset{\substack{\Vert\\\varepsilon\operatorname{int}(\beta)(\sigma)}}{\underbrace{\eta\sigma}},\\[1ex]
\det\underline{u}&=\mu
\end{aligned}\right.
\qquad\text{avec }\xi=\alpha\rho(\alpha)^{-1},\quad
\eta=\varepsilon\beta\sigma(\beta)^{-1}
\]
\margin{attention au facteur $\varepsilon$!}

\uline{Proposition} Le morphisme de schémas $\Theta$ est
\uline{caractérisé} par les formules (29), mod un morphisme
$\mathcal{G}_{\rho,\sigma}^{*}\xrightarrow{\lambda}\mu_2$.

En effet, les premières deux formules, qui décrivent l'action de
$\underline{u}$ sur $\rho$ resp. $\sigma$, déterminent $\underline{u}$ mod
$T_\rho$ resp. mod $T_\sigma$, donc \inserted{(ensemble)} mod
$T_\rho\cap T_\sigma=\mathbb{G}_{\mathrm m}$, la troisième le détermine
mod $SG$, -- donc avec les précédentes, mod $\mu_2$ -- i.e. un autre
morphisme $\Theta'$ ayant les $\overline{\mathrm m}\mathrm{s}$ propriétés
sera de la forme
\[
\Theta'(g)=\lambda(g)\Theta(g)\qquad\text{où }\lambda:
\mathcal{G}_{\rho,\sigma}^{*}\longrightarrow\unsure{\mathcal{E}_2|\varepsilon_2}.
\]
\unsure{cqfd}

\uline{Rmq} Notons \inserted{de plus} que \deleted{\unsure{tout}}
$\Theta$ est normalisé par
\[
\Theta(\underbrace{1,1,1,1}_{1})=1
\]
ce qui impose \deleted{\unsure{une}} normalisation
\deleted{\unsure{supplémentaire}}\inserted{correspondante}
$\lambda(1)=1$; sur $\lambda$ -- mais \unsure{sauf erreur} cela
implique déjà (au moins si $\rho,\sigma$ sont $\frac12$ simples donc
$T_\rho,T_\sigma$ des tores\ldots) que $\lambda$ est un \uline{hom.}
de schémas en groupes, qui ici est \unsure{nécessairement} trivial,
sur \inserted{la composante} $B'_0\times T_\rho\times T_\sigma$
(du moins en car. $\ne2$) donc on trouve essentiellement comme seule
indétermination possible la multiplication \deleted{par} de
$\underline{u}$ par $\varepsilon$ \deleted{\unsure{tout}}
(correspondant à l'hom.\inserted{can.} (non trivial) de
$N'_\sigma/T_\sigma$ sur $\mu_2$)\ldots\ \deleted{\unsure{Err.}}

\margin{Si on travaillait \unsure{sur} $S\mathcal{G}_{\rho,\sigma}$
au lieu de $\mathcal{G}_{\rho,\sigma}^{*}$, l'indétermination serait
donnée par la multiplication \unsure{par|$\mapsto$}
$\varepsilon^{i}\varepsilon'^{j}$
$(i,j\in\mathbb{Z}/2\mathbb{Z})$.}
\margin{377}

% Printed archive footer (not autograph).
\begin{center}
377\\
\texttt{https://grothendieck.umontpellier.fr}\\
Université de Montpellier, UMR CNRS 5149, Case courrier 051, Place Eugène Bataillon\\
34095 Montpellier cedex 5 - France
\end{center}

\scanpage{379}
\mspage{\unsure{490}}

(IX) \uline{Les groupes} $\mathcal{M}_{\rho,\sigma}$ \uline{et}
$S\mathcal{M}_{\rho,\sigma}$.

Soient tjrs $\rho,\sigma$ deux sections de $G$, satisfaisant les
conditions générales (cf p. ex. prop. p. \unsure{478'|470'}), plus la
condition $\operatorname{Tr}\sigma=0$ (cf \uline{V} -- p. ex. lemme
p. 455 et son cor. p. 457) -- nous pourrions faire leur construction
aussi sans cette condition supplémentaire (en nous contentant de
travailler avec $\mathcal{G}_{\rho,\sigma}$ au lieu de
$\mathcal{G}_{\rho,\sigma}^{*}$) -- mais \unsure{peu importe}
sans doute pour la suite. Nous pouvons au besoin supposer
$\rho,\sigma$ normalisés de façon qu'on ait
$\operatorname{Tr}(\rho+\sigma)=1$.

On désigne par $S\mathcal{M}_{\rho,\sigma}$ le schéma produit
\[
(1)\qquad S\mathcal{M}_{\rho,\sigma}
=S\mathcal{G}_{\rho,\sigma}^{*}\times_S SG
\]
\margin{\unsure{Par|Pour} la suite, l'``\unsure{officialisation}''
\illegible{1} des \unsure{schémas|structures} en groupes $S_0\mathcal{M}_{\rho,\sigma}$,
cf p. 559 ff.}

Ses sections sont notées
\[
(\alpha,\beta,\mu,\varepsilon,f)\ ;
\]
$f$ étant donc une section arbitraire de $SG$
$(\simeq\mathrm{Sl}(2)_S)$.

Dans le cas où $S=\operatorname{Spec}(A)$, $A=\widehat{\mathbb{Z}}$,
\[
G=\mathrm{Gl}(2)_S,\quad
\rho=\begin{pmatrix}0&1\\-1&1\end{pmatrix},\quad
\sigma=\begin{pmatrix}0&-1\\1&0\end{pmatrix},
\]
et où on a donc un homomorphisme canonique
\[
(2)\qquad
\widehat{\unsure{\mathcal{G}|\mathcal{C}}}_{0,3}
\longrightarrow
G(\widehat{\mathbb{Z}})^{\prime\pm1}
\simeq\unsure{\mathrm{Gl}|\mathrm{Sl}}(2,\widehat{\mathbb{Z}})^{\pm1}/\{\pm1\}
\]
à valeurs dans le sous-groupe des \unsure{élémts|éléments} de dét
$\pm1$,
\margin{378}

% Printed archive footer (not autograph).
\begin{center}
378\\
\texttt{https://grothendieck.umontpellier.fr}\\
Université de Montpellier, UMR CNRS 5149, Case courrier 051, Place Eugène Bataillon\\
34095 Montpellier cedex 5 - France
\end{center}

\scanpage{380}

envoyant $\widehat{\unsure{\mathcal{G}|\mathcal{C}}}_{0,3}^{+}$ dans
$SG(\widehat{\mathbb{Z}})'=\mathrm{Sl}(2,\widehat{\mathbb{Z}})'
=\mathrm{Sl}(2,\widehat{\mathbb{Z}})/\unsure{\{\pm1\}|\pm1}$,
on va construire une variante du groupe $\mathcal{M}_{0,3}^{\sim}$,
\inserted{(cf p. 390)}, qui s'envoie canoniquement dans le produit (1),
ce qui nous amènera à définir une structure de groupe sur
$S\mathcal{M}_{\rho,\sigma}$, de façon à rendre ledit morphisme
multiplicatif.

Considérons d'abord la variante ``spéciale''
$S\mathcal{M}_{0,3}^{\sim}$ de $\mathcal{M}_{0,3}^{\sim}$, extension
de $\mathcal{M}_{0,3}^{\sim}$ par $\widehat{\mathbb{Z}}^3$,
déduite d'une extension de
\[
\unsure{\mathcal{A}|\mathcal{H}}_{0,3}^{\sim}
=\mathcal{M}_{0,3}^{\sim}/\widehat{\Pi}_{0,3}
\subset\unsure{\widehat{\widehat{\mathcal{T}}}_{0,3}|\widehat{\mathcal{T}}_{0,3}},
\]
\deleted{\unsure{p}} et $\overline{\mathrm m}$ d'une extension de
$\unsure{\widehat{\widehat{\mathcal{T}}}_{0,3}|\widehat{\mathcal{T}}_{0,3}}$
(via $\mathcal{M}_{0,3}^{\sim}\longrightarrow
\unsure{\widehat{\widehat{\mathcal{T}}}_{0,3}|\widehat{\mathcal{T}}_{0,3}}$),
formée des
\[
(3)\qquad
(u,\unsure{f_0|g_0},\unsure{f_1|g_1},\unsure{f_\infty|g_\infty})
\qquad
\left(u\in\unsure{\widehat{\widehat{\mathcal{T}}}_{0,3}|\widehat{\mathcal{T}}_{0,3}},\quad
g_0,g_1,g_\infty\in\unsure{\widehat{\Pi}_{0,3}}\right)
\]
satisfaisant
\[
(4)\qquad
u(\ell_i)=\operatorname{int}\bigl(\unsure{f|g}_{\dot{u}(i)}^{-1}\bigr)
\bigl(\ell_{\dot{u}(i)}^{\mu}\bigr)
\]
où $\mu$ est le multiplicateur de $u$, et $\dot{u}$ l'action de $u$
sur \inserted{l'ens.} $\{0,1,\infty\}$ des classes de lacets.
Le produit est défini par
\[
\begin{aligned}
(5)\qquad &(u',f'_0,f'_1,f'_\infty)
(\deleted{\unsure{u}},f_0,f_1,f_\infty)\\
&=(u'u,\ f'_0u'(f_0),\ f'_1u'(f_1),\ \unsure{f'_\infty u'(f_\infty)})
\end{aligned}
\]
en nous bornant (\unsure{dégonflage!}) au cas
$u,u'\in\mathcal{M}_{0,3}^{!\sim}$ i.e. $\dot{u},\dot{u}'=1$.

\margin{379}

% Printed archive footer (not autograph).
\begin{center}
379\\
\texttt{https://grothendieck.umontpellier.fr}\\
Université de Montpellier, UMR CNRS 5149, Case courrier 051, Place Eugène Bataillon\\
34095 Montpellier cedex 5 - France
\end{center}

\scanpage{381}
\mspage{491}

On a un hom surjectif $S\mathcal{M}_{0,3}^{\sim}\longrightarrow
\mathcal{M}_{0,3}^{\sim}$, dont le noyau est formé des
\[
(6)\qquad (1,f_0,f_1,f_\infty)\qquad f_i\in L_i
\]
\unsure{donc on a}
\[
(7)\qquad
\diagram{\begin{tikzcd}[ampersand replacement=\&,column sep=large]
1\arrow[r]\&\widehat{\mathbb{Z}}^3\arrow[r]\&
S\mathcal{M}_{0,3}^{\sim}\arrow[r]\&
\mathcal{M}_{0,3}^{\sim}\arrow[r]\&1
\end{tikzcd}}
\]

Nous allons définir un sous-groupe intéressant
\[
(8)\qquad S_0\mathcal{M}_{0,3}^{\sim}\subset
\mathcal{M}_{0,3}^{\sim},
\]
en utilisant maintenant la définition de $\mathcal{M}_{0,3}^{\sim}$,
savoir \inserted{\unsure{via}} la propriété des
$u\in\mathcal{M}_{0,3}^{\sim}$ de commuter extérieurement à
$\rho$, de sorte qu'on aura
\[
(9)\qquad u\rho u^{-1}=g\rho\qquad\text{i.e. }[u,\rho]=g,
\quad g\in\widehat{\Pi}_{0,3}
\]

Ceci permet, quand on a \uline{choisi} déjà
$f_i\in\widehat{\Pi}_{0,3}$ satisfaisant
\[
(10)\qquad u(\ell_i)=\operatorname{int}(f_i^{-1})(\ell_i^{\mu})
\]
(en supposant pour simplifier $u\in\mathcal{M}_{0,3}^{!\sim}$),
d'en déduire les \deleted{\unsure{formes|formules}}
$f_{i+1},f_{i+2}$ via la formule
\[
\begin{aligned}
(11)\qquad f_{i+1}^{-1}&=g\rho(f_i^{-1})\,,\qquad
\inserted{\left[\begin{aligned}
\text{et }f_{i+2}^{-1}&=g\rho(f_{i+1}^{-1})\\
&=g\rho(g)\rho^2(f_i^{-1})
\end{aligned}\right]}\\[-1ex]
&\inserted{[\text{où }g=[u,\rho]\in\widehat{\Pi}_{0,3})}
\end{aligned}
\]
(on a posé $\dot{\rho}(i)=i+1$ par abus de langage,
i.e. $0+1=1,\ 1+1=\infty,\ \infty+1=0$!) -- on
\unsure{retrouve|retrouvera} alors
\[
f_i^{-1}=\unsure{f|g}_{(i+2)+1}^{-1}
=g\rho(f_{i+2}^{-1})
=\underbrace{g\rho(g)\rho^2(g)}_{1}\,
\underbrace{\rho^3(f_i^{-1})}_{f_i^{-1}}.
\]
\margin{380}

% Printed archive footer (not autograph).
\begin{center}
380\\
\texttt{https://grothendieck.umontpellier.fr}\\
Université de Montpellier, UMR CNRS 5149, Case courrier 051, Place Eugène Bataillon\\
34095 Montpellier cedex 5 - France
\end{center}

\scanpage{382}
\mspage{492}

Les $(u,f_0,f_1,f_\infty)
\in S\mathcal{M}_{0,3}^{!\sim}$ qui satisfont (11) pour
$\forall\ i\in\{0,1,\infty\}$ (il suffit qu'il \unsure{le|les}
vérifie pour deux indices, pour que ce soit vérifié pour tous les
trois) forment un \uline{sous-groupe} de
$S\mathcal{M}_{0,3}^{!\sim}$, que je note
$S_0\mathcal{M}_{0,3}^{!\sim}$ -- il s'identifie aussi au groupe
des couples $(u,f_0)$ $(u\in\mathcal{M},\ f_0\in\widehat{\Pi}_{0,3})$
satisfaisant
\[
(12)\qquad u(\ell_0)=\operatorname{int}(f_0^{-1})(\ell_0^{\mu})
\qquad(\inserted{\text{où }}\mu=\chi(u)).
\]

On a une structure d'extension
\[
(13)\qquad
\diagram{\begin{tikzcd}[ampersand replacement=\&,column sep=large]
1\arrow[r]\&\widehat{\mathbb{Z}}\arrow[r,no head]\&
S_0\mathcal{M}_{0,3}^{!\sim}\arrow[r]\&
\mathcal{M}_{0,3}^{!\sim}\arrow[r,no head]\&1
\end{tikzcd}},
\]
et un hom d'extensions
\[
(14)\qquad\left\{
\diagram{\begin{tikzcd}[ampersand replacement=\&,column sep=large,row sep=large]
1\arrow[r,no head]\&\widehat{\mathbb{Z}}\arrow[r,no head]\arrow[d,"\text{diagonale}"]\&
S_0\mathcal{M}_{0,3}^{!\sim}\arrow[r,no head]\arrow[d]\&
\mathcal{M}_{0,3}^{!\sim}\arrow[r,no head]\arrow[d]\arrow[d,no head,shift right=0.5ex]\&1\\
1\arrow[r,no head]\&\widehat{\mathbb{Z}}^3\arrow[r,no head]\&
S\mathcal{M}_{0,3}^{!\sim}\arrow[r,no head]\&
\mathcal{M}_{0,3}^{!\sim}\arrow[r,no head]\&1
\end{tikzcd}}\right.
\]
(où $S\mathcal{M}_{0,3}^{!\sim}$ est \inserted{défini comme}
l'image \unsure{d'envers|d'inverse} de
$\mathcal{M}_{0,3}^{!\sim}\unsure{\subset|\in}\mathcal{M}_{0,3}^{\sim}$
dans $S\mathcal{M}_{0,3}^{\sim}$). On peut donc dire que
l'extension $S\mathcal{M}_{0,3}^{!\sim}$ (de
$\mathcal{M}_{0,3}^{!\sim}$ par $\widehat{\mathbb{Z}}^3$) est
déduite, par l'hom diagonal
\[
(15)\qquad
\diagram{\begin{tikzcd}[ampersand replacement=\&,column sep=large,row sep=small]
\widehat{\mathbb{Z}}\arrow[r]\&\widehat{\mathbb{Z}}^3\\
\underline{n}\arrow[r,mapsto]\&
(\underline{n}_0=\underline{n},\ \underline{n}_1=\underline{n},\
\underline{n}_\infty=\underline{n})
\end{tikzcd}},
\]
d'une extension ``plus fine'' $S_0\mathcal{M}_{0,3}^{!\sim}$ de
\margin{381}

% Printed archive footer (not autograph).
\begin{center}
381\\
\texttt{https://grothendieck.umontpellier.fr}\\
Université de Montpellier, UMR CNRS 5149, Case courrier 051, Place Eugène Bataillon\\
34095 Montpellier cedex 5 - France
\end{center}

\scanpage{383}
$\mathcal M_{0,3}^{!\sim}$ par $\widehat{\mathbb Z}$.

Les calculs du \S 46 \uline{V} (p. 407 \unsure{ff}) permettent d'autre part d'expliciter les couples $(\unsure{u|\alpha},f_0=f)$ en termes de systèmes
\[
(16)\qquad
(\alpha,\beta,f)\quad\left|\quad
\begin{gathered}
\alpha\in\widehat{\widetilde{\Pi}}_{0,3}
 =\widehat{\Pi}_{0,3}\cdot\{1,\tau_\infty\}
 =\text{\unsure{image inv.|im. récipr.} de }\{1,\tau\}\\
\text{par }\widehat{\unsure{\mathcal G|\mathcal C}}_{0,3}\longrightarrow\mathcal T_{0,3}\\
\beta,f\in\widehat{\Pi}_{0,3}
\end{gathered}\right.
\]
\margin{à $(u,f)$ on associe $(\unsure{fu},f)$ en mettant $\unsure{fu}\in\mathcal M_{0,3}^{!\sim}[0]$ sous la forme $u_{\alpha,\beta}$}
satisfaisant des relations (ne faisant pas intervenir $f$)
\[
(17)\qquad
\alpha\rho(\alpha)^{-1}
 =\beta\sigma(\beta)^{-1}
 \underbrace{\varepsilon_1(\mu-1)}_{=\unsure{\ell_1|h}^{\nu}\ \text{si }\mu-1=2\nu}
\]
où $\mu\in\widehat{\mathbb Z}^{*}$ est défini \deleted{\unsure{sans aucune}} \unsure{ambiguité|univoquement} par $\alpha$ --- et $\overline{\mathrm m}$ \unsure{par la seule connaissance} \deleted{de l'image} de $\xi=\underline\alpha\rho(\underline\alpha)^{-1}$, où $\underline\alpha$ est l'image de $\alpha$ dans $\mathrm{Gl}(2,\widehat{\mathbb Z})^{\pm1}$.

Le cas $f=1$ correspond à \inserted{un} sous-groupe de $S_0\mathcal M_{0,3}^{!\sim}$, qui s'identifie au normalisateur de $L_0$ dans $\mathcal M_{0,3}^{!\sim}$, qu'on avait noté $\mathcal M_{0,3}^{!\sim}[0]$ (\S 46 \uline{V} $\ldots$) :
\[
\begin{gathered}
(18)\qquad
\underset{\substack{\Vert\ \mathrm{dfn}\\
\operatorname{Norm}_{\mathcal M_{0,3}^{!\sim}}(L_0)}}{\mathcal M_{0,3}^{!\sim}[0]}
\simeq
\left\{(\alpha,\beta,f)\,\middle|\,
\begin{array}{l}f=1,\ \alpha,\beta\\\text{satisfaisant }(17)\end{array}\right\}
\subset S_0\mathcal M_{0,3}^{!\sim} .
\end{gathered}
\]
Si on désigne par $\unsure{u_{\xi,\eta}|u_{\alpha,\beta}}$ \deleted{\unsure{ou $u_\xi$}} \deleted{l'autom} \unsure{ou simplement} $u_\xi$ l'autom défini en termes de
\[
(19)\qquad
\xi=\alpha\rho(\alpha)^{-1},\qquad
\eta=\beta\sigma(\beta)^{-1}
\]
\margin{382}

% Printed archive footer (not autograph).
\begin{center}
382\\
\texttt{https://grothendieck.umontpellier.fr}\\
Université de Montpellier, UMR CNRS 5149, Case courrier 051, Place Eugène Bataillon\\
34095 Montpellier cedex 5 - France
\end{center}

\scanpage{384}
\mspage{493}
par
\[
(20)\qquad u(\rho)=\xi\rho,\qquad u(\sigma)=\eta\sigma
\]
(NB $\eta$ est défini en termes de $\xi$ grâce à (17), où $\mu=\mu(\xi)$ est déterminé par $\xi$), on \unsure{aura} par construction (loc. cit.)
\[
(21)\qquad u_{\alpha,\beta,f}=\operatorname{int}(f^{-1})u_{\alpha,\beta}
\]
et au triplet $(\alpha,\beta,f)$ correspond l'élément $(u_{\alpha,\beta,f},f)$ de $S_0\mathcal M_{0,3}^{!\sim}$ :
\[
(22)\qquad (\alpha,\beta,f)\longmapsto (u_{\alpha,\beta},f)\in S_0\mathcal M_{0,3}^{!\sim} .
\]
Ayant ainsi explicité les éléments de $S_0\mathcal M_{0,3}^{!\sim}$, on définit une application
\[
(23)\qquad
\diagram{\begin{tikzcd}[ampersand replacement=\&,column sep=large,row sep=large]
S_0\mathcal M_{0,3}^{!\sim}\arrow[r] \&
S\mathcal M_{\rho,\sigma}(\widehat{\mathbb Z})'
=S\mathcal G_{\rho,\sigma}^{*}(\widehat{\mathbb Z})'
\times\mathrm{Sl}(2,\widehat{\mathbb Z})\unsure{'}\\
(\alpha,\beta,f)\arrow[r,mapsto] \&
\left(\underline\alpha,\underline\beta,
\mu=\mu(\underline\xi=\underline\alpha\rho(\underline\alpha)^{-1}),
\underset{\substack{\uparrow\\\inserted{\unsure{\varepsilon}}}}{\varepsilon_2(\mu)},\underline f\right)
\end{tikzcd}}
\]
où les soulignés $\underline\alpha,\underline\beta,\underline f$ \unsure{\&} désignent les images de $\alpha,\beta,f$ dans $\mathrm{Sl}(2,\widehat{\mathbb Z})'$, et l'\unsure{``accent''} $'$ dans $S\mathcal M_{\rho,\sigma}(\widehat{\mathbb Z})'$ désigne l'ens. des \deleted{couples} \inserted{systèmes} $(\underline\alpha,\underline\beta,\mu,\varepsilon,\underline f)$
\deleted{\unsure{définis d'éléments de}}
\[
\in\mathrm{Gl}(2,\widehat{\mathbb Z})^{\prime\unsure{\pm1}}
\times\mathrm{Sl}(2,\widehat{\mathbb Z})'
\times\widehat{\mathbb Z}^{*}\times\{\pm1\}
\times\inserted{\mathrm{Sl}(2,\widehat{\mathbb Z})'}
\]
satisfaisant \inserted{\uline{dans $\mathrm{Sl}(2,\widehat{\mathbb Z})$} (\uline{pas dans $\unsure{\mathrm{Sl}}(2,\widehat{\mathbb Z})'$ !}),} la relation des lacets \deleted{\unsure{pris|par} modulo passage à}
\[
(24)\qquad
\underline\alpha\rho(\underline\alpha)^{-1}
 =\varepsilon\underline\beta\sigma(\underline\beta)^{-1}\varepsilon_1(\mu-1)
\]
(ce qui ne dépend pas des choix des \uline{signes} $\ldots$)
\margin{383}

% Printed archive footer (not autograph).
\begin{center}
383\\
\texttt{https://grothendieck.umontpellier.fr}\\
Université de Montpellier, UMR CNRS 5149, Case courrier 051, Place Eugène Bataillon\\
34095 Montpellier cedex 5 - France
\end{center}

\scanpage{385}
On a ``fait ce qu'il fallait'' pour que la restriction de (23) à $\deleted{\unsure{S_0}}\mathcal M_{0,3}^{!\sim}[0]$ ((18), identifié aux $(\alpha,\beta,f)$ pour lesquels $f=1$) soit un homomorphisme \inserted{dans le groupe $S\mathcal G_{\rho,\sigma}$} --- en \unsure{définissant} justement de façon astucieuse \inserted{ad hoc} la loi de composition dans $S\mathcal G_{\rho,\sigma}^{*}$.
Je vais la revérifier quand même pour \unsure{me rassurer}, en considérant $u=u_{\alpha,\beta}$, $u'=u_{\alpha',\beta'}$, $u''=uu'$ sous la forme $u_{\alpha'',\beta''}$, et en se rappelant la ``relation des cocycles'' qui donne, à un des deux \unsure{signes} près, $(\alpha'',\beta'')$ en termes de $(\alpha',\beta'),(\alpha,\beta)$ (\S 46 \uline{II} $5^\circ$))
\[
(25)\qquad
\alpha''=u'_{\unsure{\infty|\omega}}(\alpha)\alpha',\qquad
\beta''=u'(\beta)\beta',\qquad
\inserted{\text{en plus de }\mu''=\mu\mu'}
\]
Soient \unsure{alors|donc} $\underline\alpha,\underline\beta,\underline\alpha',\underline\beta'$ des \deleted{images} \inserted{\unsure{élémts}} de $\mathrm{Gl}(2,\widehat{\mathbb Z})^{\pm1}$ (déterminés au \uline{signe} près) qui relèvent les images de $\alpha,\alpha',\beta,\beta'$ dans $\mathrm{Gl}(2,\widehat{\mathbb Z})'$, on aura donc $(\underline\alpha,\underline\beta,\mu,\varepsilon_2(\mu))$ et $(\underline\alpha',\underline\beta',\mu',\varepsilon_2(\mu'))$ dans $S\mathcal G_{\rho,\sigma}^{*}(\widehat{\mathbb Z})$, et appliquant $\Theta$ on trouve $\underline u,\underline u'\in\mathrm{Gl}(2,\widehat{\mathbb Z})$ (en fait, $\in B'_0(\widehat{\mathbb Z})$) tels que
\[
(26)\qquad
\left\{\begin{aligned}
\underline u'(\rho)&=\operatorname{int}(\underline\alpha')\rho
&\underline u'(\sigma)&=\varepsilon_2(\mu')\operatorname{int}(\underline\beta')(\sigma)\\
\det\underline u'&=\mu'
\end{aligned}\right.
\]
(et relations \unsure{item|idem} pour $u$)
et on aura alors, pour tout $x\in\deleted{\unsure{\Pi}}\widehat{\unsure{\mathcal G|\mathcal C}}_{0,3}$
\[
(27)\qquad \underline{u'(x)}=\underline u'(\underline x)
\]
(cf $\ldots$) Donc, pour \deleted{des} relèvements convenables $\underline\alpha'',\underline\beta^{\unsure{''|'}}$ des images de $\alpha^{\unsure{''|'}},\beta^{\unsure{''|'}}$ dans $\mathrm{Gl}(2,\widehat{\mathbb Z})^{\pm1}$, on aura
\[
\begin{gathered}
(28)\qquad
\underline\alpha''=\underline u'(\underline\alpha)\underline\alpha',\qquad
\underline\beta''=\underline u'(\underline\beta)\underline\beta',\\
\text{en plus de }\mu''=\mu\mu'\quad\text{d'où}\quad
\varepsilon_2''=\varepsilon_2\varepsilon_2'.
\end{gathered}
\]
\margin{384}

% Printed archive footer (not autograph).
\begin{center}
384\\
\texttt{https://grothendieck.umontpellier.fr}\\
Université de Montpellier, UMR CNRS 5149, Case courrier 051, Place Eugène Bataillon\\
34095 Montpellier cedex 5 - France
\end{center}

\scanpage{386}
\mspage{494}
et il reste à vérifier que l'élément
\[
\begin{gathered}
g''=(\underline\alpha'',\underline\beta'',\mu''=\mu\mu',\varepsilon_2''=\varepsilon_2'\varepsilon_2)
\in\deleted{\unsure{\Gamma}(\mathrm{Gl}(2,\widehat{\mathbb Z})\times\mathrm{Sl}(2,\widehat{\mathbb Z})\times}\\
\inserted{\Gamma\bigl(\widehat{\mathbb Z};\mathrm{Gl}(2)\times\mathrm{Sl}(2)\times\unsure{\mathbf E}'\times\mu_2\bigr)}
\end{gathered}
\]
est bien le produit des éléments
\[
g'=(\underline\alpha',\underline\beta',\mu',\varepsilon')\quad\text{et}\quad
g=(\underline\alpha,\underline\beta,\mu,\varepsilon).
\]
Que ce soit OK pour $\mu''$ et $\varepsilon''$ est trivial, il reste à vérifier que dans $\mathcal G_{\rho,\sigma}^{*}$, le produit est bien (en ce qui concerne les composantes en $\alpha$ et en $\beta$) donné par les formules (28) --- ce qui \unsure{sera} valable, bien entendu, sur une base quelconque, et dans les conditions générales de \uline{V} (ou \inserted{$\overline{\mathrm m}$} de \uline{III}, si on s'intéresse à $\mathcal G_{\rho,\sigma}$, au lieu de $\mathcal G_{\rho,\sigma}^{*}$). On aura
\[
(29)\qquad
\left\{\begin{aligned}
\underline\alpha&=\underline u\,\underline a^{-1},&
\underline\beta&=\underline u\,\underline b^{-1}\\
\underline\alpha'&=\underline u'\,\underline a'^{-1},&
\underline\beta'&=\underline u'\,\underline b'^{-1}
\end{aligned}\right.
\qquad
\begin{gathered}
\underline\alpha,\underline\alpha'\in T_\rho(\widehat{\mathbb Z})\\
\underline\beta,\underline\beta'\in N'_\sigma(\widehat{\mathbb Z})
\end{gathered}
\]
et par définition des produits dans $\mathcal G_{\rho,\sigma}^{*}$, le produit de \deleted{\unsure{calculé à}} $g'g$ aura donc composantes $\underline\alpha_0''$ et $\underline\beta_0''$
\[
\underline\alpha_0''=(\underline u'\underline u)(\underline a'\underline a)^{-1},\qquad
\underline\beta_0''=(\underline u'\underline u)(\underline b'\underline b)^{-1}
\]
et il reste à prouver
\[
\underline\alpha_0''=\underline\alpha'',\qquad
\underline\beta_0''=\underline\beta'',\quad\text{i.e. les relations}
\]
\[
(30)\qquad
\underbrace{\underline u'(\underline\alpha)\cdot\underline\alpha'}
 =(\underline u'\underline u)(\underline a'\underline a)^{-1}
\qquad
\underline u'(\underline\beta)\underline\beta'
 =(\underline u'\underline u)(\underline b'\underline b)^{-1}
\]
où $\deleted{\unsure{\underline u}}\ \underline\alpha,\underline\beta,\underline\alpha',\underline\beta'$ sont donnés en fonction de \deleted{$\unsure{\underline\alpha,\underline\beta,\underline u,\underline a}$} $\underline u,\underline a,\underline b,\underline u',\underline a',\underline b'$ par (29). On a donc
\[
\begin{aligned}
\underline u'(\underline\alpha)\underline\alpha'
 &=\underbrace{\underline u'(\underline u\,\underline a^{-1})\underline u'^{-1}}\underline u'\underline a'^{-1}\\
 &=(\underline u'\underline u)(\underline a^{-1}\underline a'^{-1})
 =(\underline u'\underline u)(\underline a'\underline a)^{-1},
 \qquad\deleted{\text{\unsure{OK}}}
\end{aligned}
\]
\deleted{$\unsure{\underline u'(\underline\beta)\underline\beta'}$} et de $\overline{\mathrm m}$ pour la deuxième relation 30, ce qui prouve notre assertion.
\margin{$\unsure{\checkmark|/}$}
\margin{385}

% Printed archive footer (not autograph).
\begin{center}
385\\
\texttt{https://grothendieck.umontpellier.fr}\\
Université de Montpellier, UMR CNRS 5149, Case courrier 051, Place Eugène Bataillon\\
34095 Montpellier cedex 5 - France
\end{center}

\scanpage{387}
Notons à cette occasion que les systèmes
\[
(31)\qquad
(\alpha,\beta,\unsure{\mu}=1,\varepsilon=1)
\quad\text{avec}\quad
\alpha\in\Gamma\deleted{\unsure{\{\pm1\}}}\inserted{\mathbb G_m},\qquad
\beta\in\Gamma\deleted{\unsure{\{\pm1\}}}\inserted{\mathbb G_m}
\]
sont dans $\unsure{\Gamma|S}\mathcal G_{\rho,\sigma}^{\unsure{\circ|0}}$.
\deleted{et $\overline{\mathrm m}$ dans $S\mathcal G_{\rho,\sigma}$ si $\alpha,\beta\in\Gamma\mu_2$ ---}
[qui sont stables en effet par multiplication, \deleted{\unsure{pour}} $(\alpha,\beta,\mu,\varepsilon)$, de $\alpha$ et $\beta$ chacun par une section de $\deleted{\unsure{\mu_2}}\inserted{\mathbb G_m}$, ce qui ne change pas $\widetilde\alpha,\widetilde\beta\ \ldots$]
ils correspondent, par l'isomorphisme $\mathcal G_{\rho,\sigma}\simeq B'_0\times T_\rho\times T_\sigma$, aux systèmes $(\underline u,\underline a,\underline b)$ avec $\underline u=1$, $\underline a\in\Gamma\mathbb G_m$, $\underline b\in\Gamma\mathbb G_m$) et forment un \uline{sous-groupe} (avec la multiplication ordinaire \deleted{\unsure{sur les|sur le}} \deleted{en $\alpha,\beta$} \inserted{en $\alpha\ \beta$} comme loi de groupe) --- si $\alpha,\beta$ sont sections de $\mu_2$, on trouve un \unsure{ss-groupe} \deleted{trivial} ``\unsure{de nature|d'action} triviale'' \inserted{$\mu_2\times\mu_2$} de $S\mathcal G_{\rho,\sigma}$.
L'action par translation de ce sous-groupe $\simeq\mathbb G_m\times\mathbb G_m$ sur $\mathcal G_{\rho,\sigma}^{*}$ se fait par multiplication \deleted{dans} \inserted{des} \deleted{les} composantes en question. Cela montre que le fait de travailler avec des $\underline\alpha,\underline\beta$ définis seulement ``à \uline{signes} près'' revient simplement à travailler dans les groupes quotients de \uline{$\mathcal G_{\rho,\sigma}^{*}(\widehat{\mathbb Z})$} (ou de $S\mathcal G_{\rho,\sigma}^{*}(\widehat{\mathbb Z})$) par le sous-groupe
\[
\inserted{\boxed{\begin{gathered}\Sigma_0\ .\\ \simeq\mathbb Z/2\mathbb Z\times\mathbb Z/2\mathbb Z\end{gathered}}}
\]
\deleted{engendré} des éléments $(\alpha,\beta,1,1)$ avec $\alpha,\beta\in\{\pm1\}$. Ceci
\margin{386}

% Printed archive footer (not autograph).
\begin{center}
386\\
\texttt{https://grothendieck.umontpellier.fr}\\
Université de Montpellier, UMR CNRS 5149, Case courrier 051, Place Eugène Bataillon\\
34095 Montpellier cedex 5 - France
\end{center}

\scanpage{388}
\mspage{495}

précise en $\overline{\mathrm m}$ temps le sens de la locution :
on a un \unsure{hom.} de groupes \ldots .

Revenons maintenant à l'application (23) avec $f$ quelconque, et précisons la loi de
composition de $S_0\mathcal M_{0,3}^{!\sim}$ en termes de
$(\alpha,\beta,f)$ i.e. en termes des $u_{\alpha,\beta,f}$.
Considérons le sous-groupe formé des $(1,1,f)$,
\deleted{ou encore} \deleted{des} qui correspond aux \inserted{couples}
$(u,f)$ avec $u=\operatorname{int}(f^{-1})$, i.e. aux couples
$(\operatorname{int}(f^{-1}),f)$
--- ils forment un \unsure{\textquotedblleft}groupe isomorphe à $\widehat\Pi_{0,3}$,
et on a une structure de produit $\frac12$ direct :
\[
(31)\qquad
S_0\mathcal M_{0,3}^{!\sim}
\simeq \mathcal M_{0,3}^{!\sim}[0]\cdot\widehat\Pi_{0,3}
\]
où
\[
(32)\qquad
\left\{
\diagram{\begin{tikzcd}[ampersand replacement=\&,column sep=large,row sep=small]
\widehat\Pi_{0,3} \arrow[r] \&
S_0\mathcal M_{0,3}^{!\sim} \& {}\\
f \arrow[r,mapsto] \&
(\deleted{\unsure{f,}}\operatorname{int}(f),f)
\arrow[r,leftrightarrow] \& (1,1,f^{-1}),
\end{tikzcd}}
\right.
\]
l'opération de $\mathcal M_{0,3}^{!\sim}[0]$ sur $\widehat\Pi_{0,3}$
\deleted{\unsure{étant}} l'opération tautologique
(déduite de l'inclusion
$\mathcal M_{0,3}^{!\sim}[0]\subset\mathcal M_{0,3}^{!\sim}
\subset\underline{\operatorname{Aut}}_{\mathrm{lac}}
(\widehat\Pi_{0,3})$).

Le sous-groupe \inserted{invariant} $\widehat{\mathbb Z}$ de
$S_0\mathcal M_{0,3}^{!\sim}$ est donné par les couples
$\inserted{(u,f)=}(1,\ell_0^n)$ $(n\in\widehat{\mathbb Z})$
qu'on met aussi sous la forme $(\alpha,\beta,f)$
où $\alpha,\beta$ sont tels que
$u_{\alpha,\beta}=fu=\operatorname{int}(\ell_0^n)$,
donc $\alpha=\beta=\ell_0^n$, i.e.
\[
(33)\qquad
\widehat{\mathbb Z}\longrightarrow S_0\mathcal M_{0,3}^{!\sim}
\quad\text{s'identifie à}\quad
n\longmapsto
(\deleted{\unsure{\alpha=}}\ell_0^n,\beta=\ell_0^n,f=\ell_0^n)
\]
\margin{387}

% Printed archive footer (not autograph).
\begin{center}
387\\
\texttt{https://grothendieck.umontpellier.fr}\\
Université de Montpellier, UMR CNRS 5149, Case courrier 051, Place Eugène Bataillon\\
34095 Montpellier cedex 5 - France
\end{center}

\scanpage{389}

On \unsure{aura}
\[
\begin{aligned}
(f'^{-1}u_{\alpha',\beta'})(f^{-1}u_{\alpha,\beta})
&=\deleted{\unsure{f'^{-1}f^{-1}(f u_{\alpha',\beta'}f^{-1})}}\\
&\inserted{\underbrace{f'^{-1}\unsure{u_{\alpha',\beta'}}(f^{-1})}
_{\unsure{f''^{-1}}}
(u_{\alpha',\beta'}u_{\alpha,\beta})}
\end{aligned}
\]
\unsure{i.e.} la multiplication de $S_0\mathcal M_{0,3}^{!\sim}$
se fait, en termes des composantes
$(\alpha',\beta',f')(\alpha,\beta,f)$,
\unsure{pour|par} \deleted{\unsure{g''=}}$(\alpha'',\beta'',f'')$, où
\[
(34)\quad
\left\{
\begin{aligned}
u_{\alpha'',\beta''}
&=u_{\alpha',\beta'}u_{\alpha,\beta}
\quad\text{i.e.}\quad
\left\{
\begin{aligned}
\alpha''&=u_{\alpha',\beta'}(\alpha)\alpha'\\
\beta''&=u_{\alpha',\beta'}(\beta)\beta'
\end{aligned}
\right.,
\quad\text{et}\\
f''&=u_{\alpha',\beta'}(f)f'
\end{aligned}
\right.
\]

Revenant maintenant à la définition de $S\mathcal M_{\rho\sigma}$
(1), nous allons y introduire la multiplication \unsure{correspondante},
ce qui revient à définir $S\mathcal M_{\rho\sigma}$ comme
un produit $\frac12$ direct
\[
\begin{aligned}
(35)\qquad S\mathcal M_{\rho,\sigma}
&\overset{\mathrm{dfn}}{=}
S\mathcal G_{\rho,\sigma}^{*}\cdot SG\\
&\qquad(\text{produit }\tfrac12\text{ direct})
\end{aligned}
\]
\uline{par l'opération de} $S\mathcal G_{\rho,\sigma}^{*}$
\uline{sur} $G$
\uline{via l'homomorphisme} $\Theta$ (p.~\unsure{489'|488'}),
et \unsure{en} \unsure{définissant|définira}
\unsure{le|l'} \deleted{application}\inserted{morphisme}
\[
(36)\qquad
\left|
\diagram{\begin{tikzcd}[ampersand replacement=\&,column sep=large,row sep=small]
S\mathcal M_{\rho,\sigma}\arrow[r] \&
S\mathcal G_{\rho,\sigma}^{*}\times_S G\\
\deleted{\unsure{(\alpha,\beta,\mu,\varepsilon)|(\alpha,\beta)}}
\inserted{U}\cdot f\arrow[r,mapsto] \& (U,f^{-1})
\end{tikzcd}}
\right.
\]
\margin{Attention au passage de $f\underset{\leftarrow}{\ \unsure{\text{à}}\ }f^{-1}$ !}
où $U=(\alpha,\beta,\mu,\varepsilon)$ est dans
$\Gamma S\mathcal G_{\rho,\sigma}$.
\margin{388}

% Printed archive footer (not autograph).
\begin{center}
388\\
\texttt{https://grothendieck.umontpellier.fr}\\
Université de Montpellier, UMR CNRS 5149, Case courrier 051, Place Eugène Bataillon\\
34095 Montpellier cedex 5 - France
\end{center}

\scanpage{390}
\mspage{496}

On a donc fait ce qu'il fallait, dans le cas
$S=\operatorname{Spec}(\widehat{\mathbb Z})$ etc, pour avoir un
\uline{homomorphisme de groupes}
\[
(37)\qquad
\diagram{\begin{tikzcd}[ampersand replacement=\&,column sep=large]
S_0\mathcal M_{0,3}^{!\sim}\arrow[r] \&
S\mathcal M_{\rho,\sigma}(\widehat{\mathbb Z})/
\unsure{\widetilde{\Sigma}}\deleted{\unsure{(\widehat{\mathbb Z})}}
\end{tikzcd}}
\]
où $\Sigma_0\simeq\mathbb Z_{/2}\times\mathbb Z_{/2}$
est le \unsure{ss-groupe} \deleted{\unsure{\textquotedblleft impuissant\textquotedblright}}
de $S\mathcal G_{\rho,\sigma}^{*}$ introduit p.~494',
qui opère trivialement sur $G$ donc est invariant dans
$S\mathcal M_{\rho,\sigma}$, donc dans
$S\mathcal M_{\rho,\sigma}(\widehat{\mathbb Z})$, et où
\[
(38)\quad
\widetilde{\Sigma}=\Sigma_0\times
\underbrace{\{\pm1\}}_{
\substack{
\text{ce facteur identifié au centre de }SG(\unsure{\widehat{\mathbb Z}})\\
=\operatorname{Sl}(2,\widehat{\mathbb Z}),\
\text{contenu dans }\operatorname{Gl}(2,\widehat{\mathbb Z})\\
=G(\widehat{\mathbb Z})
}}
\simeq
\mathbb Z/2\mathbb Z\times\mathbb Z/2\mathbb Z\times\mathbb Z/2\mathbb Z
\]
est le groupe distingué correspondant dans le produit $\frac12$ direct
--- le fait de passer au quotient par $\Sigma$ exprime
les trois \unsure{ambiguïtés|ambiguités} de \uline{signes},
sur $\alpha$, sur $\beta$ et sur $f$.

Considérons maintenant l'\unsure{hom}
\[
(39)\quad
\left|
\diagram{\begin{tikzcd}[ampersand replacement=\&,column sep=large,row sep=small]
\mathbb G_a\arrow[r] \&
\unsure{S_0|S}\mathcal M_{\rho,\sigma}^{\deleted{\unsure{!\sim}}}\\
\lambda\arrow[r,mapsto] \&
(\alpha\unsure{=}\varepsilon_0(\unsure{\lambda|-\lambda}),
\beta\unsure{=}\varepsilon_0(-\lambda),
\mu=1,\varepsilon=1,f=\varepsilon_0(\lambda))
\end{tikzcd}}
\right.
\]
c'est un hom. de groupes, et fait $\overline{\mathrm m}$ de
$\mathbb G_a$ un ss-groupe \uline{distingué}.
Pour \inserted{vérifier} ceci le plus commode est d'identifier
$\unsure{S_0|S}\mathcal M_{\rho,\sigma}^{!\sim}$
au produit $\frac12$-direct (35),
\margin{389}

% Printed archive footer (not autograph).
\begin{center}
389\\
\texttt{https://grothendieck.umontpellier.fr}\\
Université de Montpellier, UMR CNRS 5149, Case courrier 051, Place Eugène Bataillon\\
34095 Montpellier cedex 5 - France
\end{center}

\scanpage{391}

où on identifie $S\mathcal G_{\rho,\sigma}^{*}$ à un sous-groupe
de $B'_0\times T_\rho\times N'_\sigma$, opérant sur $G$
par autom. int. \uline{via} $B'_0$.

Le morphisme précédent correspond donc à
\[
(40)\qquad
\diagram{\begin{tikzcd}[ampersand replacement=\&,column sep=small,row sep=small]
\lambda\arrow[r,mapsto,"i"]\arrow[d,phantom,"\in"{rotate=90}] \&
{(\lambda,1,1,} \&
\lambda^{-1}\arrow[d,phantom,"\in"{rotate=90}] \&
{)=\deleted{\lambda\cdot\unsure{j}(\lambda)}}\\
\Gamma B'_0 \& {} \& \Gamma G \& {}
\end{tikzcd}}
\]
\margin{\deleted{\diagram{\begin{tikzcd}[ampersand replacement=\&,column sep=small]
\unsure{j}:SG\arrow[r] \& S\mathcal M_{\rho,\sigma}
\end{tikzcd}} est} l'hom de groupes canonique}

\margin{\deleted{\unsure{où $U_0$}}}
\deleted{Soit d'autre part
$U=(\underbrace{\underline u,\underline{\unsure{a|\rho}},
\underline{\unsure{b|\beta}}}_{U_0},\unsure{\varepsilon}f)$
\inserted{$=U_0\cdot f$},
un élément quelconque du produit $\frac12$ direct,
et calculons $U i(\lambda)U^{-1}$ ---
\[
U i(\lambda)U^{-1}=U_0 f i(\lambda)f^{-1}U_0^{-1}=
\]
}

On identifie $S\mathcal G_{\rho,\sigma}^{*}$ à un sous-groupe de
$\mathcal G_{\rho,\sigma}^{*}$, et de $\overline{\mathrm m}$
le produit $\frac12$ direct $S\mathcal M_{\rho,\sigma}^{\unsure{*}}$
à un sous-groupe du produit $\frac12$ direct
\[
(41)\qquad
\deleted{\unsure{S\mathcal M_{\rho,\sigma}^{*}}}
\inserted{\widetilde S\mathcal M_{\rho,\sigma}}
=\mathcal G_{\rho,\sigma}^{*}\cdot G
\simeq(B'_0\times T_\rho\times N'_\sigma)\cdot G,
\]
et on va voir que l'image du \unsure{mono} $i$ (40) dans
$\mathcal G_{\rho,\sigma}^{*}\cdot G$ est $\overline{\mathrm m}$
distinguée dans ce groupe plus grand.
Pour ne pas nous y perdre dans les notations,
désignons par
\[
(42)\qquad i_B,\ i_\rho,\ i_\sigma,\ i_G
\]
les morphismes \unsure{d'inclusions|d'inclusion}
\[
\left|
\diagram{\begin{tikzcd}[ampersand replacement=\&,column sep=large,row sep=small,nodes in empty cells]
B'_0\arrow[r] \& \deleted{\unsure{\mathcal G_{\rho,\sigma}^{*}}}\\
T_\rho\arrow[r] \& {}\\
\unsure{N'_\sigma|N''_\sigma}\arrow[r] \& {}\\
G\arrow[r] \& {}
\end{tikzcd}}
\right\}\ 
\unsure{i'}
\begin{gathered}
\inserted{\widetilde S\mathcal M_{\rho,\sigma}}\\
\deleted{\unsure{\mathcal G_{\rho,\sigma}^{*}|\mathcal M_{\rho,\sigma}}}\cdot G
\end{gathered}
\]
\margin{390}

% Printed archive footer (not autograph).
\begin{center}
390\\
\texttt{https://grothendieck.umontpellier.fr}\\
Université de Montpellier, UMR CNRS 5149, Case courrier 051, Place Eugène Bataillon\\
34095 Montpellier cedex 5 - France
\end{center}

\scanpage{392}
\mspage{497}

étant entendu, bien sûr, qu'on fait opérer
$\mathcal G_{\rho,\sigma}^{*}$ sur $G$ intérieurement via le
facteur $B'_0$. Il faut vérifier que
\[
i_{\deleted{\unsure{B}}}(L_0)
=\{i_B(\lambda)i_G(\lambda^{-1})\mid\lambda\in\Gamma L_0\}
\]
est invariant par
$\deleted{\unsure{\Gamma T_\rho},\unsure{\Gamma T_\sigma}\unsure{i_T}}$
\inserted{\unsure{les}}
$i_\rho(x)$, $i_\sigma(x)$, $i_B(x)$ et $i_G(x)$.
--- Pour les deux premiers, c'est trivial --- avec action triviale sur $L_0$.

Si $x\in\Gamma B'_0$, on aura
\[
\begin{aligned}
i_B(x)(i(\lambda))
&=i_B(x)(i_B(\lambda)i_G(\lambda^{-1}))=\\
&=\underbrace{i_B(x)(i_B(\lambda))}
_{i_B(\deleted{\unsure{\lambda|x}}x(\lambda))}
\cdot
\underbrace{i_B(x)(i_G(\lambda^{-1}))}
_{\substack{
i_G(x)(i_G(\lambda^{-1}))\\
=i_G(x(\lambda)^{-1}]
}}\\
&=i(x(\lambda))
\end{aligned}
\]
\margin{On \unsure{voit} de suite que $i(\lambda)$
commute à $G$, donc est invariant par $G$ !}
ce qui montre que $\deleted{\unsure{\mathbb G_a|L_0}}$
\[
\diagram{\begin{tikzcd}[ampersand replacement=\&,column sep=large]
L_0\arrow[r,"i"] \&
\mathcal G_{\rho,\sigma}^{*}\cdot G
=\widetilde S\mathcal G_{\rho,\sigma}^{*}
\end{tikzcd}}
\]
commute à l'action naturelle de $B'_0$ sur les deux côtés
(sur le deuxième via $\unsure{\pi_B}$,
\[
\diagram{\begin{tikzcd}[ampersand replacement=\&,column sep=large]
i_B:B'_0\arrow[r] \&
\widetilde S\underset{\deleted{\unsure{0}}}{\mathcal M}_{\rho,\sigma}^{\unsure{\flat|*|\natural}}
\end{tikzcd}}
\]
). Enfin
\[
i_G(\deleted{\unsure{f}}\inserted{\unsure{x}})(i(\lambda))
=i_G(x)\cdot(i_B(\lambda)i_G(\lambda^{-1}))\cdot i_G(x)^{-1}
\]
or par définition, $i_B(\lambda)$ et $i_G(\lambda)$ opèrent
de la $\overline{\mathrm m}$ façon sur $G$, donc
$i(\lambda)=i_B(\lambda)i_G(\lambda^{-1})$
y opère trivialement, donc le dernier facteur
$i_G(x)^{-1}$ peut \unsure{sauter} \unsure{sur|sous} le facteur
$i(\lambda)$, on trouve
\[
i_G(x)i_G(x)^{-1}i(\lambda)=i(\lambda)
\]
\margin{O.K.}
\margin{391}

% Printed archive footer (not autograph).
\begin{center}
391\\
\texttt{https://grothendieck.umontpellier.fr}\\
Université de Montpellier, UMR CNRS 5149, Case courrier 051, Place Eugène Bataillon\\
34095 Montpellier cedex 5 - France
\end{center}

\scanpage{393}

\uline{Rmq} $\widetilde S\mathcal M_{\rho,\sigma}$ est un groupe plat de dim relative
$\underbrace{2}_{B'_0}+\underbrace{2}_{T_\rho}+\underbrace{2}_{N'_\sigma}+\underbrace{4}_{G}=10$,
$S\mathcal M_{\rho,\sigma}$ est de dim relative
$\underbrace{\unsure{4}}_{S\mathcal G_{\rho,\sigma}^{*}}+\underbrace{3}_{SG}=7$,
en passant au quotient par $\underset{\simeq\mathbb G_a}{i(L_0)}$
\unsure{fournit|fournir} des groupes \deleted{\unsure{ce}} notés
$\widetilde{\mathcal M}_{\rho,\sigma}$, $\mathcal M_{\rho,\sigma}$,
de dim relative $9$ resp. $6$, donnant lieu à un diagramme
\[
(43)\qquad
\diagram{\begin{tikzcd}[ampersand replacement=\&,column sep=large,row sep=large]
1 \arrow[r] \& \mathbb G_a \arrow[r,"\lambda\mapsto\varepsilon_0(\lambda)"] \arrow[d,equal]
  \& S\mathcal M_{\rho,\sigma} \arrow[r] \arrow[d,hook,"\unsure{\mathrm{can.}}"']
  \& \mathcal M_{\rho,\sigma} \arrow[r] \arrow[d,hook] \& 1 \\
1 \arrow[r] \& \mathbb G_a \arrow[r]
  \& \widetilde S\mathcal M_{\rho,\sigma} \arrow[r]
  \& \unsure{\mathcal M_{\rho,\sigma}|\widetilde{\mathcal M}_{\rho,\sigma}} \arrow[r,no head] \& 1
\end{tikzcd}}
\]

Le passage au quotient ne fait pas de difficultés --- les quotients sont en fait
\unsure{représentables} \deleted{\unsure{par les|par}}\inserted{comme} produits $1/2$ directs
\[
(44)\qquad
\left\{\begin{aligned}
\deleted{\unsure{S}}\mathcal M_{\rho,\sigma}
 &\simeq \deleted{\unsure{T_0}}\,S\mathcal G_{\rho,\sigma}^{*}(0)\cdot SG\\
\widetilde{\mathcal M}_{\rho,\sigma}^{\deleted{\unsure{\sim}}}
 &\simeq \mathcal G_{\rho,\sigma}^{*\deleted{\unsure{\flat|\natural|\circ}}}(\unsure{0|\omega})\cdot\deleted{\unsure{S}}G
\end{aligned}\right.
\]
où les sous-schémas en groupes
$S\mathcal G_{\rho,\sigma}^{*\unsure{\flat|\natural|\circ}}(\unsure{0|\omega})$ et
$\mathcal G_{\rho,\sigma}^{*\unsure{\flat|\natural|\circ}}(\unsure{0|\omega})$ de
$S\mathcal G_{\rho,\sigma}^{*}$, $\mathcal G_{\rho,\sigma}^{*}$
sont définis dans $\boxed{\text{\uline{VIII}}}$,
$\mathcal G_{\rho,\sigma}^{*\unsure{\flat|\natural|\circ}}(\unsure{0|\omega})$ s'identifiant à
$T_0\times T_\rho\times N'_\sigma$, et
$S\mathcal G_{\rho,\sigma}^{*\unsure{\flat|\natural|\circ}}(0)$ à son intersection
avec $S\mathcal G_{\rho,\sigma}^{*}$ (défini par les conditions suppl. $(20)$ de p. $487'$).

\margin{392}

% Printed archive footer (not autograph).
\begin{center}
392\\
\texttt{https://grothendieck.umontpellier.fr}\\
Université de Montpellier, UMR CNRS 5149, Case courrier 051, Place Eugène Bataillon\\
34095 Montpellier cedex 5 - France
\end{center}

\scanpage{394}
\mspage{498}

Revenant maintenant à la situation arithmétique, on trouve
\inserted{par passage au quotient via $(37)$} un \unsure{hom.} canonique
\margin{compte tenu de $(13)$}
\[
(45)\qquad
\diagram{\begin{tikzcd}[ampersand replacement=\&,column sep=large]
\deleted{\unsure{S_0}}\mathcal M_{0,3}^{!\sim} \arrow[r]
\& \underbrace{\mathcal M_{\rho,\sigma}(\widehat{\mathbb Z})}_{\begin{gathered}
\simeq\\
\underbrace{\deleted{\unsure{B'_0}}\unsure{T_0|\widehat T_0}\times
T_\rho(\widehat{\mathbb Z})\times T_\sigma(\widehat{\mathbb Z})}_{\begin{gathered}
\text{opérant sur}\\
\mathrm{Sl}(2,\widehat{\mathbb Z})\ \text{via le}\\
\text{facteur }T_0
\end{gathered}}\cdot\mathrm{Sl}(2,\widehat{\mathbb Z})
\end{gathered}}/\Sigma
\end{tikzcd}}
\]
$\Sigma$ étant le sous-groupe $\simeq\{\pm1\}^3$, produit
de trois facteurs $\mathbb Z/2\mathbb Z$ inclus resp.$^{\unsure{t}}$
dans $T_\rho(\widehat{\mathbb Z})$, $T_\sigma(\widehat{\mathbb Z})$
et dans $\mathrm{Sl}(2,\widehat{\mathbb Z})$.

\uline{Rmq} Cet homomorphisme a été obtenu ici,
\inserted{essentiellement,} à l'aide de l'hom. de
$\mathcal M_{0,3}^{!\sim}\unsure{(0)}$
\[
(46)\qquad
\diagram{\begin{tikzcd}[ampersand replacement=\&,column sep=large,row sep=huge]
\begin{gathered}
\mathcal M_{0,3}^{!\sim}(0)\\
\text{sous-groupe}\\
\text{des autom.}\\
\text{``normalisés''}
\end{gathered} \arrow[dr]
\& \begin{gathered}
=\operatorname{Norm}_{\mathcal M_{0,3}^{!\sim}}(L_0)\\
\text{(cf }\S46,\ \text{\uline{V}}\text{)}
\end{gathered} \arrow[dr,-{Circle[length=2pt,width=2pt]}] \& \\
\& S\mathcal G_{\rho,\sigma}^{*\unsure{\flat|\natural|\circ}}(0)(\widehat{\mathbb Z})/\Sigma_0
  \arrow[r,hook,no head]
\& S\mathcal G_{\rho,\sigma}^{*}(\widehat{\mathbb Z})/\Sigma_0
\end{tikzcd}}
\]
et du composé
\[
\diagram{\begin{tikzcd}[ampersand replacement=\&,column sep=large,row sep=small]
\deleted{\unsure{\mathcal G}}\widehat\Pi_{0,3}
  \arrow[r,hook,no head]
\& \widehat{\unsure{\mathcal G|\mathcal C}}_{0,3}^{+}
  \arrow[r] \arrow[d,phantom,"\unsure{\cup|\wr}" description]
\& \mathrm{Sl}(2,\widehat{\mathbb Z})' \\
\& \mathrm{Sl}(2,\mathbb Z)^{\prime\wedge} \&
\end{tikzcd}}
\]
en écrivant
\[
(47)\qquad
\mathcal M_{0,3}^{!\sim}\simeq\mathcal M_{0,3}^{!\sim}\cdot\widehat\Pi_{0,3}\quad\unsure{!|:}
\]

\margin{393}

% Printed archive footer (not autograph).
\begin{center}
393\\
\texttt{https://grothendieck.umontpellier.fr}\\
Université de Montpellier, UMR CNRS 5149, Case courrier 051, Place Eugène Bataillon\\
34095 Montpellier cedex 5 - France
\end{center}

\scanpage{395}

Mais on a aussi un produit semi direct
\[
(48)\qquad
\mathcal M_{0,3}^{\sim}\simeq
\mathcal M_{0,3}^{!\sim}(\unsure{0|\omega})\cdot
\unsure{\mathcal G|\mathcal C}_{0,3}^{+}
\]
qui, combiné directement à
\[
\diagram{\begin{tikzcd}[ampersand replacement=\&,column sep=large]
\widehat{\unsure{\mathcal G|\mathcal C}}_{0,3}^{+}\arrow[r]
\&\mathrm{Sl}(2,\widehat{\mathbb Z})'
\end{tikzcd}}
\]
donne un hom
\[
(49)\qquad
\diagram{\begin{tikzcd}[ampersand replacement=\&,column sep=large]
\mathcal M_{0,3}^{\sim}\arrow[r]
\& \bigl(\mathcal M_{\rho,\sigma}(\widehat{\mathbb Z})
\simeq S\mathcal G_{\rho,\sigma}^{*}(\widehat{\mathbb Z})\cdot
\mathrm{Sl}(2,\widehat{\mathbb Z})\bigr)/\unsure{\Sigma}
\end{tikzcd}}
\]
qui prolonge l'hom $(45)$, défini d'abord sur son sous-groupe
distingué $\mathcal M_{0,3}^{!\sim}$ d'indice $6$.

\margin{394}

% Printed archive footer (not autograph).
\begin{center}
394\\
\texttt{https://grothendieck.umontpellier.fr}\\
Université de Montpellier, UMR CNRS 5149, Case courrier 051, Place Eugène Bataillon\\
34095 Montpellier cedex 5 - France
\end{center}

\scanpage{396}
\mspage{499}

$\overline{\text{\uline{X}}}$\quad Étude du groupe
$\mathcal M_{\rho,\sigma}$ et de $\mathcal M_{\rho,\sigma}\to\mathrm{Gl}(2)$

Par construction du produit semi-direct
\[
(1)\qquad
\diagram{\begin{tikzcd}[ampersand replacement=\&,column sep=small,row sep=huge]
\mathcal M_{\rho,\sigma}^{\sim} \arrow[r,phantom,"=" description]
\& \mathcal G_{\rho,\sigma}^{*\unsure{\flat|\natural|\circ}}(\unsure{\omega|0})\cdot G
  \arrow[r,phantom,"\simeq" description] \arrow[d,phantom,"\cup" description]
\& \underbrace{(T_0\times T_\rho\times N'_\sigma)}_{\begin{gathered}
\text{opérant sur }G\\ \text{via }T_0\hookrightarrow G
\end{gathered}}\cdot G \\
\mathcal M_{\rho,\sigma} \arrow[r,phantom,"=" description]
\& S\mathcal G_{\rho,\sigma}^{*\unsure{\flat|\natural|\circ}}(\unsure{\omega|0})\cdot SG \arrow[r,hook]
\& \underbrace{(T_\rho\times N'_\sigma)}_{\begin{gathered}
\text{opérant sur }G\\
\text{via }N'_\sigma\xrightarrow{\det}\mathbb G_m
\xleftarrow[\sim]{\det}T_0\hookrightarrow G
\end{gathered}}\cdot G
\end{tikzcd}}
\]
on a un homomorphisme de groupes
\[
(2)\qquad
\left\{\begin{gathered}
\diagram{\begin{tikzcd}[ampersand replacement=\&,column sep=large]
\mathcal M_{\rho,\sigma}^{\sim}\arrow[r,"\Theta"] \& G
\end{tikzcd}}\\
\Theta\bigl(i_B(x)i_{\deleted{\unsure{\rho|B}}\inserted{\rho}}(y)
\unsure{i_\sigma|i'_\sigma}(z)i_G(f)\bigr)=x.f
\end{gathered}\right.
\]
\[
(x\in\Gamma T_0,\quad y\in\Gamma T_\rho,\quad
z\in\Gamma\unsure{N'_\sigma},\quad f\in\Gamma G).
\]

L'homomorphisme composé
\[
(3)\qquad
\diagram{\begin{tikzcd}[ampersand replacement=\&,column sep=large,row sep=large]
S\mathcal M_{\rho,\sigma}^{\sim}\arrow[r,no head]
  \arrow[d,phantom,"\unsure{\wr}" description]
\&\mathcal M_{\rho,\sigma}^{\sim}\arrow[r,"\Theta"] \& G\\
\underbrace{\deleted{\unsure{\mathcal M}}\mathcal G_{\rho,\sigma}^{*\unsure{\flat|\natural|\circ}}}_{B'_0\times T_\rho\times N'_\sigma}\cdot G
\&\&
\end{tikzcd}}
\]
\margin{(composé noté encore $\Theta$)}
\deleted{est} décrit par
\[
\Theta\bigl(i_B(x)i_\rho(y)i_\sigma(z)i_G(\unsure{f|\varphi})\bigr)=x.f
\]
et $\Theta$ sera bien trivial pour les
\[
i(\lambda)=i_B(\lambda)i_G(\lambda^{-1})\qquad(\lambda\in\Gamma L_0)
\]
et passe donc au quotient en $\Theta:\mathcal M_{\rho,\sigma}^{\sim}\to G$.

\margin{395}

% Printed archive footer (not autograph).
\begin{center}
395\\
\texttt{https://grothendieck.umontpellier.fr}\\
Université de Montpellier, UMR CNRS 5149, Case courrier 051, Place Eugène Bataillon\\
34095 Montpellier cedex 5 - France
\end{center}
